\subsection*{Infraestructura como código frente a infraestructura dinámica}
\addcontentsline{toc}{subsection}{Infraestructura como código frente a infraestructura dinámica}

El concepto de \textit{Infrastructure as Code} (IaC) se basa en aplicar prácticas propias del desarrollo de software a la gestión de sistemas e infraestructura. Esta metodología propone definir el entorno tecnológico mediante código, permitiendo que los procesos de provisión, configuración y actualización de los sistemas sean automáticos, repetibles y verificables.

Según Morris (2016) \cite{Morris2016}, la clave de IaC es tratar la infraestructura como si fuera código fuente, lo que hace posible aprovechar herramientas de control de versiones, pruebas automatizadas y orquestación de despliegues en su mantenimiento. De esta forma, se integran prácticas consolidadas del desarrollo ágil —como la Integración Continua (CI) y la Entrega Continua (CD)— en la administración de entornos en la nube.

Antes de adoptar el modelo de \textit{Infrastructure as Code} (IaC), es importante comprender las limitaciones que presenta la llamada infraestructura dinámica, un enfoque común en las primeras etapas de adopción del cómputo en la nube. Este modelo se basa en la provisión y configuración manual o semiautomatizada de servidores y entornos, normalmente mediante interfaces gráficas o scripts ad hoc. Aunque inicialmente resulta flexible y rápido para desplegar nuevos recursos, a medida que el sistema crece aparecen problemas de mantenimiento, consistencia y escalabilidad que afectan de forma directa la fiabilidad del entorno.

Uno de los principales riesgos es la proliferación de servidores y configuraciones difíciles de gestionar, fenómeno que Morris (2016) \cite{Morris2016} denomina “\textit{server sprawl}”. Cuando el equipo técnico puede crear recursos con facilidad pero sin control centralizado, el número de servidores aumenta más rápido que la capacidad de administrarlos correctamente. Esto genera entornos heterogéneos, con versiones de software y configuraciones distintas, donde los errores se multiplican y las tareas de actualización o parcheo se vuelven complejas y arriesgadas.

A esta diversidad no controlada se le suma la llamada "deriva de configuración" (\textit{configuration drift}), que ocurre cuando servidores idénticos en su origen se modifican con el tiempo de forma independiente. Pequeños ajustes, correcciones manuales o actualizaciones parciales provocan diferencias sutiles que terminan impidiendo reproducir el entorno completo, generando lo que Morris describe como “servidores copo de nieve” (\textit{snowflake servers}): máquinas únicas e irrepetibles que nadie se atreve a modificar por miedo a romper su funcionamiento. Cuando esta fragilidad se extiende a todo el ecosistema, se obtiene una “infraestructura Jenga”, donde cualquier cambio puede desestabilizar el conjunto.

En paralelo, la falta de estandarización genera lo que el autor denomina “miedo a la automatización” (\textit{automation fear}). Los equipos evitan ejecutar herramientas de automatización de manera continua, debido a la falta de confianza en los resultados o al temor de introducir inconsistencias en sistemas ya frágiles. Como consecuencia, las configuraciones se aplican de forma irregular y la automatización se limita a tareas puntuales, perpetuando un ciclo de desconfianza y errores.

Frente a estos problemas, el modelo de \textit{Infrastructure as Code} ofrece una alternativa estructurada y sostenible. Mediante la definición declarativa y versionada de la infraestructura, es posible reconstruir cualquier entorno desde cero de forma repetible, auditable y libre de dependencias ocultas. Cada cambio queda registrado en el control de versiones, integrándose con los flujos de integración y despliegue continuo. Este enfoque elimina la dependencia de configuraciones manuales, reduce la deriva entre entornos y evita la formación de servidores especiales no replicables.

Por estas razones, en el proyecto \textbf{Muncher} se opta por un modelo basado en IaC. Ello permite mantener coherencia entre entornos, asegurar la reproducibilidad del sistema y facilitar la automatización completa del proceso de despliegue. Además, proporciona una base sólida para escalar la infraestructura, incorporar nuevas funcionalidades y responder con agilidad ante cambios o incidencias en la nube.
